%  ★ 자작 2단 조판본(USV_swarm_defense_EAAI_KR.tex) 전용 사본이다.
%    EAAI 투고본은 elsarticle 의 highlights 환경(sections/00c_highlights.tex)을 쓴다.
%============================================================ 연구 하이라이트
%  ⚠ "3~5 항목, 85자 이내" 는 로컬 Elsevier 배포물에서 확인되지 않는 통념이다.
%  ★ 길이는 tools/check_paper_meta.py 가 기계 검사한다 — 눈대중으로 세지 말 것.
%  각 항목은 "무엇을 했다"가 아니라 "무엇이 밝혀졌다"로 쓴다.
%  영문판은 참고문헌 뒤 English Abstract 절에 함께 둔다(표제 span 과적재 방지).

\vspace{4pt}
\begin{center}
\begin{minipage}{0.96\textwidth}
\small
\noindent\textbf{연구 하이라이트}
\begin{itemize}[leftmargin=1.2em, itemsep=1pt, topsep=2pt, parsep=0pt]
\item 같은 25초 주기의 두 판단 계층을 지연 차이만으로 비동기 분리한 방어 체계
\item 점수맵 픽셀 지목으로 행동공간을 격자에 가둔, 적선 수에 불변인 협동 정책
\item 가치망 없는 그룹 상대 최적화와 반사실적 크레딧으로 다중 방어정 협동 학습
\item 720 에피소드 대응표본 분해 결과 두 계층의 기여는 10지표 중 9지표에서 가산적
\item 지휘관 성능은 배정 커버리지가 아니라 재계획 간 계획 안정성이 설명한다
\end{itemize}
\end{minipage}
\end{center}
